(1) 在球内外无自由体电荷且介电常数各自均匀时，有 $\nabla^2\phi=0$。对轴对称且远场只有 $l=1$ 分量的解，可写
\[
\phi_{\rm in}=-AE_0r\cos\theta,\qquad
\phi_{\rm out}=-E_0r\cos\theta+\frac{B\cos\theta}{r^2},
\]
其中球心处已舍去 $1/r^2$ 奇异项，外部的额外增长项也由远场条件排除。$A$ 无量纲，$B$ 具有电势乘长度平方的量纲。

(2) 在 $r=a$，电势连续给 $-AE_0a=-E_0a+B/a^2$，即 $B=E_0a^3(1-A)$。法向 $D$ 连续：球内 $E_{{\rm in},r}=AE_0\cos\theta$，球外 $E_{{\rm out},r}=E_0\cos\theta+2B\cos\theta/a^3$，故
\[
\epsilon A E_0=\epsilon_0(E_0+2B/a^3).
\]
消元得
\[
A=\frac{3\epsilon_0}{\epsilon+2\epsilon_0},\qquad
B=\frac{\epsilon-\epsilon_0}{\epsilon+2\epsilon_0}E_0a^3.
\]

(3) 介质球内电场沿 $+\hat z$ 均匀：
\[
\boldsymbol E_{\rm in}=\frac{3\epsilon_0}{\epsilon+2\epsilon_0}\boldsymbol E_0,\qquad
\boldsymbol P=(\epsilon-\epsilon_0)\boldsymbol E_{\rm in}.
\]
体束缚电荷 $-\nabla\cdot\boldsymbol P=0$，表面束缚电荷
\[
\sigma_b(\theta)=\boldsymbol P\cdot\hat{\boldsymbol r}
=\frac{3\epsilon_0(\epsilon-\epsilon_0)}{\epsilon+2\epsilon_0}E_0\cos\theta.
\]
球面上 $\int_0^\pi\cos\theta\sin\theta\,d\theta=0$，因此总束缚电荷为零。将球外 $B\cos\theta/r^2$ 与真空偶极子电势 $p\cos\theta/(4\pi\epsilon_0r^2)$ 对比，得到
\[
\boldsymbol p=4\pi\epsilon_0a^3\frac{\epsilon-\epsilon_0}{\epsilon+2\epsilon_0}\boldsymbol E_0.
\]

(4) 当 $\epsilon\to\epsilon_0$ 时 $A\to1,B\to0,\boldsymbol P\to0$，球变得与真空不可区分；当 $\epsilon\to\infty$ 时 $A\to0$，球内电场趋零，而
\[
\boldsymbol p\to4\pi\epsilon_0a^3\boldsymbol E_0,
\]
恰为理想导体球在均匀外静电场中的感应偶极矩。不能把极限 $\epsilon\to\infty$ 下的极化强度误判为零：虽然 $E_{\rm in}\to0$，但 $(\epsilon-\epsilon_0)E_{\rm in}$ 趋向有限非零值。
